\chapter{Quản lý dữ liệu}

\section{Nguồn dữ liệu và giấy phép}

Dữ liệu của dự án gồm ba nhóm. Thứ nhất là audio/transcript do người dùng upload khi demo, được xem là dữ liệu nhạy cảm và không commit vào Git. Thứ hai là worker roster do người dùng cung cấp, dùng để chuẩn hoá assignee. Thứ ba là synthetic benchmark \texttt{data/eval/gold\_synthetic\_205.jsonl}, gồm 205 meeting samples với transcript turns, participants và action items tham chiếu. Synthetic benchmark được tạo cho kiểm thử nội bộ nên không phụ thuộc giấy phép dữ liệu bên thứ ba.

\section{Kích thước dữ liệu và ngôn ngữ}

Benchmark chính có 205 meetings. Khi huấn luyện baseline action-turn detector, các meetings được tách thành 2036 transcript turns. Metadata checkpoint ghi nhận 1620 train turns, 416 test turns và positive rate 0.3001. Benchmark hiện ưu tiên tiếng Anh, phù hợp với prompt/model runtime đang dùng trong demo.

\section{Tiền xử lý và làm sạch dữ liệu}

Audio runtime được chuẩn hoá trước khi đưa vào ASR. Transcript được tổ chức thành \codefield{TranscriptTurn} gồm speaker, start/end timestamp và text. Với benchmark synthetic, script baseline gán nhãn turn bằng cách so khớp token overlap giữa action item reference và từng transcript turn. Dữ liệu được validate schema, kiểm tra record rỗng và giữ meeting-level grouping để tránh leakage.

\section{Lược đồ dữ liệu}

Thực thể trung tâm của hệ thống là \texttt{Meeting}. Mỗi meeting sinh ra transcript turns, action items, participant mapping, feedback corrections, artifacts và calendar events. Thiết kế này giữ được đường truy vết từ task quay lại đoạn hội thoại gốc, đồng thời tách dữ liệu review/feedback khỏi output ban đầu của model.

Hình \ref{fig:data-schema} mô tả quan hệ giữa meeting, transcript, task, participant, worker roster, feedback và artifact. Bảng \ref{tab:data-schema} mô tả schema trung gian quan trọng nhất được pipeline tạo ra hoặc tiêu thụ.

\begin{figure}[H]
\centering
\resizebox{\textwidth}{!}{%
\begin{tikzpicture}[
  entity/.style={draw=primary, fill=light, rounded corners=3pt, text width=3.2cm, minimum height=0.8cm, align=center, font=\scriptsize},
  arrow/.style={-Latex, thick, color=darkgray}
]
  \node[entity, fill=stagecolor1] (meeting) {\textbf{meetings}\\id, status, summary, model\_version};
  \node[entity, above left=1.0cm and 1.6cm of meeting] (turns) {\textbf{transcript\_turns}\\turn\_id, speaker, time, text};
  \node[entity, above right=1.0cm and 1.6cm of meeting] (tasks) {\textbf{tasks}\\description, assignee, due, source turns};
  \node[entity, left=2.0cm of meeting] (participants) {\textbf{meeting\_participants}\\speaker\_id, display\_name, worker\_id};
  \node[entity, right=2.0cm of meeting] (feedback) {\textbf{feedback\_corrections}\\original/corrected fields};
  \node[entity, below left=1.0cm and 1.6cm of meeting] (artifacts) {\textbf{meeting\_artifacts}\\payload, checksum, metadata};
  \node[entity, below right=1.0cm and 1.6cm of meeting] (calendar) {\textbf{calendar\_events}\\provider, event id, status};
  \node[entity, below=2.1cm of meeting, fill=stagecolor2] (workers) {\textbf{workers}\\worker\_id, name, email, aliases};

  \draw[arrow] (meeting) -- (turns);
  \draw[arrow] (meeting) -- (tasks);
  \draw[arrow] (meeting) -- (participants);
  \draw[arrow] (meeting) -- (feedback);
  \draw[arrow] (meeting) -- (artifacts);
  \draw[arrow] (meeting) -- (calendar);
  \draw[arrow, dashed] (participants) -- (workers);
  \draw[arrow, dashed] (tasks) -- (workers);
\end{tikzpicture}
}
\caption{Lược đồ dữ liệu chính của hệ thống}
\label{fig:data-schema}
\end{figure}

\begin{table}[H]
\centering
\footnotesize
\caption{Schema dữ liệu trung gian và output chính}
\label{tab:data-schema}
\begin{tabularx}{\textwidth}{@{}>{\raggedright\arraybackslash}p{3.9cm}X@{}}
\toprule
\textbf{Đối tượng} & \textbf{Trường chính} \\
\midrule
\codefield{TranscriptTurn} & \codefield{turn_id}, \codefield{speaker_id}, \codefield{speaker_name}, \codefield{start_ms}, \codefield{end_ms}, \texttt{text}, \codefield{asr_confidence}. \\
\codefield{ExtractedTask} & \codefield{task_id}, \texttt{description}, \texttt{assignee}, \codefield{assignee_id}, \codefield{due_date}, \texttt{priority}, \codefield{source_turn_ids}, \texttt{status}. \\
\codefield{MeetingSummary} & \codefield{meeting_id}, \codefield{job_status}, \texttt{participants}, \codefield{transcript_turns}, \codefield{summary_text}, \codefield{action_items}, \codefield{run_metrics}, \codefield{model_version}. \\
\codefield{FeedbackCorrection} & \codefield{task_id}, \texttt{original\_*}, \texttt{corrected\_*}, \codefield{is_false_positive}, \codefield{is_missing}, \codefield{submitted_at}. \\
\bottomrule
\end{tabularx}
\end{table}

\section{Chia tập train/test và chất lượng dữ liệu}

Baseline checkpoint dùng meeting-level split 80/20: các turns từ cùng một meeting không xuất hiện đồng thời ở train và test. Cách chia này phản ánh đúng bản chất dữ liệu cuộc họp, vì các lượt nói trong cùng meeting có ngữ cảnh phụ thuộc nhau.

Dữ liệu synthetic hỗ trợ reproducibility và regression testing. Khi mở rộng sang dữ liệu thật, pipeline feedback sẽ cung cấp corrections từ người dùng như missing task, false positive, corrected assignee và corrected due date. Các correction này được export JSONL để phục vụ evaluation hoặc training sau khi scrub PII, validate schema và loại trùng.

\section{Dữ liệu thiếu, nhiễu và thiên lệch}

Hệ thống xử lý dữ liệu thiếu hoặc mơ hồ bằng trạng thái review thay vì lưu output như kết quả đã xác nhận. Audio nhiễu, speaker overlap hoặc transcript thiếu ngữ cảnh được phản ánh qua job status, guardrail signals và human feedback. Về bias, ASR/diarization có thể hoạt động khác nhau theo accent, chất lượng microphone, ngôn ngữ hoặc speaking style; vì vậy benchmark tương lai cần bao phủ nhiều dạng meeting hơn.
